\documentclass[tikz,border=6pt]{standalone}
\usepackage{ctex}
\usepackage{amsmath,amssymb}
\usetikzlibrary{arrows.meta,positioning,fit,calc,backgrounds}

\begin{document}

\begin{tikzpicture}[
    font=\small,
    >=Latex,
    node distance=8mm and 10mm,
    block/.style={
        draw,
        rectangle,
        rounded corners=1.5pt,
        align=center,
        minimum height=10mm,
        text width=38mm,
        inner sep=4pt
    },
    eqblock/.style={
        draw,
        rectangle,
        rounded corners=1.5pt,
        align=center,
        minimum height=12mm,
        text width=46mm,
        inner sep=5pt
    },
    smallblock/.style={
        draw,
        rectangle,
        rounded corners=1.5pt,
        align=center,
        minimum height=9mm,
        text width=28mm,
        inner sep=3pt
    },
    line/.style={->, thick},
    groupbox/.style={
        draw,
        rounded corners=2pt,
        inner sep=6pt
    },
    grouplabel/.style={
        anchor=south west,
        fill=white,
        inner xsep=3pt,
        inner ysep=1pt,
        font=\small
    }
]

% Top inputs
\node[block] (cmd) {%
导引指令\\
$u_d,\ \psi_d$
};

\node[block, below left=of cmd] (eu) {%
速度误差\\
$e_u = u_d-u$
};

\node[block, below right=of cmd] (epsi) {%
航向误差\\
$e_\psi = \operatorname{wrap}(\psi_d-\psi)$
};

% PID equations
\node[eqblock, below=of eu] (pidu) {%
纵向 PID 控制律\\[1mm]
$\tau_u = K_{pu}e_u + K_{iu}\!\int e_u\,dt + K_{du}\dot e_u$
};

\node[eqblock, below=of epsi] (pidpsi) {%
航向 PID 控制律\\[1mm]
$\tau_r = K_{p\psi}e_\psi + K_{i\psi}\!\int e_\psi\,dt + K_{d\psi}\dot e_\psi$
};

% Output combination
\node[eqblock, below=16mm of $(pidu)!0.5!(pidpsi)$] (tau) {%
控制输入合成\\[1mm]
$\tau=\begin{bmatrix}\tau_u\\[1mm]0\\[1mm]\tau_r\end{bmatrix}$
};

% Optional allocation
\node[block, below=of tau] (alloc) {%
若为双推进器 USV：\\
$\tau \mapsto [T_L,\ T_R]^\top$
};

% States
\node[smallblock, left=18mm of eu] (state1) {%
状态反馈\\
$u$
};

\node[smallblock, right=18mm of epsi] (state2) {%
状态反馈\\
$\psi$
};

% connections
\draw[line] (cmd) -- (eu);
\draw[line] (cmd) -- (epsi);

\draw[line] (eu) -- (pidu);
\draw[line] (epsi) -- (pidpsi);

\draw[line] (pidu) |- (tau);
\draw[line] (pidpsi) |- (tau);

\draw[line] (tau) -- (alloc);

\draw[line] (state1) -- (eu);
\draw[line] (state2) -- (epsi);

% groups
\begin{scope}[on background layer]
    \node[groupbox, fit=(cmd)] (g1) {};
    \node[grouplabel] at ([xshift=2mm,yshift=1mm]g1.north west) {指令输入};

    \node[groupbox, fit=(eu)(epsi)(state1)(state2)] (g2) {};
    \node[grouplabel] at ([xshift=2mm,yshift=1mm]g2.north west) {误差构造};

    \node[groupbox, fit=(pidu)(pidpsi)] (g3) {};
    \node[grouplabel] at ([xshift=2mm,yshift=1mm]g3.north west) {PID 控制律};

    \node[groupbox, fit=(tau)(alloc)] (g4) {};
    \node[grouplabel] at ([xshift=2mm,yshift=1mm]g4.north west) {控制输出};
\end{scope}

\end{tikzpicture}

\end{document}